% =========================================================
% Master Summary: Balls, Open Sets, Closed Sets
% Source: volume-iii/analysis/metric-spaces/notes/open-closed-sets/
%
% Purpose: A single reference page collecting EVERY way a ball
%          or set can be characterised as open or closed.
%          No new material — pure synthesis.
%
% Dependencies (all results proved elsewhere in the chapter):
%   def:open-ball, def:closed-ball            (notes-balls.tex)
%   def:open-closed-set                       (open-closed-sets.tex)
%   def:interior, def:exterior, def:boundary  (open-closed-sets.tex)
%   def:closure, def:adherent-point           (open-closed-sets.tex)
%   def:limit-point, def:derived-set          (notes-limit-points.tex)
%   prop:balls-open-closed                    (open-closed-sets.tex)
%   prop:open-as-union-closed                 (notes-working-with-balls.tex)
%   prop:closed-as-int-open                   (notes-working-with-balls.tex)
%   thm:ibe-decomp                            (open-closed-sets.tex)
%   thm:boundary-closure                      (open-closed-sets.tex)
%   prop:closed-owns-boundary                 (open-closed-sets.tex)
%   thm:closed-equiv                          (notes-sequential-characterizations.tex)
%   prop:sequential-duality                   (notes-sequential-characterizations.tex)
%   thm:closure-three-way-unifying            (notes-sequential-characterizations.tex)
% =========================================================

% ---------------------------------------------------------
\subsubsection{Master Reference: Balls, Open Sets, Closed Sets}
\label{sec:open-closed-master-summary}
% ---------------------------------------------------------

\begin{tcolorbox}[colback=gray!6, colframe=gray!40, arc=2pt,
  left=8pt, right=8pt, top=6pt, bottom=6pt,
  title={\small\textbf{How to Use This Page}},
  fonttitle=\small\bfseries]
This page collects every characterisation of balls and open/closed sets
from across the chapter into one place.
Nothing here is new --- every row is proved elsewhere and
linked back to its source.
Use it to check at a glance which test applies to your proof,
or to see how the various perspectives connect.
\end{tcolorbox}

\bigskip

% =========================================================
% PART I: BALLS
% =========================================================

\begin{tcolorbox}[colback=gray!6, colframe=gray!40, arc=2pt,
  left=8pt, right=8pt, top=6pt, bottom=6pt,
  title={\small\textbf{Part I \enspace Balls}},
  fonttitle=\small\bfseries]
\small
\renewcommand{\arraystretch}{1.55}
\begin{tabular}{p{3.8cm} p{6.2cm} p{3.6cm}}
\toprule
\textbf{Object} & \textbf{Characterisation} & \textbf{Source} \\
\midrule

\textbf{Open ball} $B_r(x)$
  & $B_r(x) = \{y \in X : d(x,y) < r\}$
  & \hyperref[def:open-ball]{D\ref*{def:open-ball}} \\[2pt]

\textbf{Closed ball} $\overline{B}_r(x)$
  & $\overline{B}_r(x) = \{y \in X : d(x,y) \le r\}$
  & \hyperref[def:closed-ball]{D\ref*{def:closed-ball}} \\[6pt]

$B_r(x)$ \textbf{is open}
  & Radial gap: $y \in B_r(x) \Rightarrow B_\delta(y) \subseteq B_r(x)$
    where $\delta = r - d(x,y) > 0$
  & \hyperref[prop:balls-open-closed]{P\ref*{prop:balls-open-closed}} \\[2pt]

$\overline{B}_r(x)$ \textbf{is closed}
  & Radial gap: $y \notin \overline{B}_r(x) \Rightarrow B_\delta(y) \cap \overline{B}_r(x) = \varnothing$
    where $\delta = d(x,y) - r > 0$
  & \hyperref[prop:balls-open-closed]{P\ref*{prop:balls-open-closed}} \\[6pt]

\textbf{Open ball from closed}
  & $B_r(x) = \bigcup_{0 < s < r} \overline{B}_s(x)$
  & \hyperref[prop:open-as-union-closed]{P\ref*{prop:open-as-union-closed}} \\[2pt]

\textbf{Closed ball from open}
  & $\overline{B}_r(x) = \bigcap_{\varepsilon>0} B_{r+\varepsilon}(x)$
  & \hyperref[prop:closed-as-int-open]{P\ref*{prop:closed-as-int-open}} \\[6pt]

\textbf{Warning: }$B_r(x) \subsetneq \overline{B}_r(x)$ \textbf{in general}
  & In the discrete metric: $B_1(x) = \{x\}$ but
    $\overline{B}_1(x) = X$. The two balls can be very different.
  & Example \\[2pt]

\bottomrule
\end{tabular}
\end{tcolorbox}

\bigskip

% =========================================================
% PART II: OPEN SETS — ALL CHARACTERISATIONS
% =========================================================

\begin{tcolorbox}[colback=gray!6, colframe=gray!40, arc=2pt,
  left=8pt, right=8pt, top=6pt, bottom=6pt,
  title={\small\textbf{Part II \enspace Open Sets — Every Characterisation}},
  fonttitle=\small\bfseries]
\small
Let $(X,d)$ be a metric space and $U \subseteq X$.
All of the following are equivalent.
\renewcommand{\arraystretch}{1.55}
\begin{tabular}{p{2.0cm} p{6.4cm} p{3.8cm}}
\toprule
\textbf{Language} & \textbf{$U$ is open iff \ldots} & \textbf{Source} \\
\midrule

\textbf{Ball (def)}
  & $\forall x \in U,\; \exists r > 0 : B_r(x) \subseteq U$
  & \hyperref[def:open-closed-set]{D\ref*{def:open-closed-set}} \\[2pt]

\textbf{Interior}
  & $U = \operatorname{int}(U)$
    \newline (every point of $U$ is an interior point)
  & \hyperref[def:interior]{D\ref*{def:interior}} \\[2pt]

\textbf{Boundary}
  & $U \cap \partial U = \varnothing$
    \newline (open sets contain none of their boundary)
  & \hyperref[prop:closed-owns-boundary]{P\ref*{prop:closed-owns-boundary}} \\[2pt]

\textbf{Complement}
  & $X \setminus U$ is closed
  & \hyperref[def:open-closed-set]{D\ref*{def:open-closed-set}} \\[2pt]

\textbf{Sequential}
  & $x_n \to x \in U \Rightarrow x_n \in U$ eventually
    \newline (approaching sequences eventually enter $U$)
  & \hyperref[prop:sequential-duality]{P\ref*{prop:sequential-duality}} \\[2pt]

\textbf{IBE decomp}
  & $U = \operatorname{int}(U) \sqcup \varnothing$:
    every point of $U$ lies in $\operatorname{int}(U)$,
    none in $\partial U$ or $\operatorname{ext}(U)$
  & \hyperref[thm:ibe-decomp]{T\ref*{thm:ibe-decomp}} \\[2pt]

\bottomrule
\end{tabular}

\medskip
\noindent\textbf{One-line summary:} \;
$U$ open $\iff$ every point of $U$ has room around it $\iff$
$U$ contains no boundary points $\iff$ limits approaching $U$ eventually enter $U$.
\end{tcolorbox}

\bigskip

% =========================================================
% PART III: CLOSED SETS — ALL CHARACTERISATIONS
% =========================================================

\begin{tcolorbox}[colback=gray!6, colframe=gray!40, arc=2pt,
  left=8pt, right=8pt, top=6pt, bottom=6pt,
  title={\small\textbf{Part III \enspace Closed Sets — Every Characterisation}},
  fonttitle=\small\bfseries]
\small
Let $(X,d)$ be a metric space and $F \subseteq X$.
All of the following are equivalent.
\renewcommand{\arraystretch}{1.55}
\begin{tabular}{p{2.0cm} p{6.4cm} p{3.8cm}}
\toprule
\textbf{Language} & \textbf{$F$ is closed iff \ldots} & \textbf{Source} \\
\midrule

\textbf{Complement (def)}
  & $X \setminus F$ is open
  & \hyperref[def:open-closed-set]{D\ref*{def:open-closed-set}} \\[2pt]

\textbf{Sequences}
  & $(x_n) \subseteq F,\; x_n \to x \Rightarrow x \in F$
    \newline (limits of sequences in $F$ stay in $F$)
  & \hyperref[thm:closed-equiv]{T\ref*{thm:closed-equiv}} \\[2pt]

\textbf{Limit points}
  & $F' \subseteq F$
    \newline ($F$ contains all its limit points)
  & \hyperref[thm:closed-equiv]{T\ref*{thm:closed-equiv}} \\[2pt]

\textbf{Closure}
  & $F = \overline{F}$
    \newline ($F$ equals its own closure)
  & \hyperref[def:closure]{D\ref*{def:closure}} \\[2pt]

\textbf{Boundary}
  & $\partial F \subseteq F$
    \newline ($F$ contains all its boundary points)
  & \hyperref[prop:closed-owns-boundary]{P\ref*{prop:closed-owns-boundary}} \\[2pt]

\textbf{IBE decomp}
  & $\operatorname{ext}(F) = X \setminus F$:
    the exterior of $F$ is exactly the complement,
    so $\partial F$ falls inside $F$
  & \hyperref[thm:ibe-decomp]{T\ref*{thm:ibe-decomp}} \\[2pt]

\bottomrule
\end{tabular}

\medskip
\noindent\textbf{One-line summary:} \;
$F$ closed $\iff$ $F$ absorbs all its limits $\iff$
$F$ contains all its limit points $\iff$
$F$ owns its boundary $\iff$ $F$ equals its closure.
\end{tcolorbox}

\bigskip

% =========================================================
% PART IV: CLOSURE — THREE LANGUAGES
% =========================================================

\begin{tcolorbox}[colback=gray!6, colframe=gray!40, arc=2pt,
  left=8pt, right=8pt, top=6pt, bottom=6pt,
  title={\small\textbf{Part IV \enspace Closure — Three Languages}},
  fonttitle=\small\bfseries]
\small
Let $(X,d)$ be a metric space, $E \subseteq X$, $x \in X$.
The following are equivalent \textnormal{(\hyperref[thm:closure-three-way-unifying]{Unifying Theorem})}.
\renewcommand{\arraystretch}{1.55}
\begin{tabular}{p{2.2cm} p{4.8cm} p{5.4cm}}
\toprule
\textbf{Language} & \textbf{Statement} & \textbf{Meaning} \\
\midrule

\textbf{(T) Topological}
  & $x \in \overline{E}$
  & $x$ is an adherent point of $E$ \\[2pt]

\textbf{(G) Geometric}
  & $\forall r > 0,\; B_r(x) \cap E \neq \varnothing$
  & every ball around $x$ hits $E$ \\[2pt]

\textbf{(S) Sequential}
  & $\exists\, (a_n) \subseteq E : a_n \to x$
  & $x$ is a limit of a sequence in $E$ \\[2pt]

\bottomrule
\end{tabular}

\medskip
\begin{tabular}{p{6cm} p{6.5cm}}
\toprule
\textbf{What closure adds to $E$} & \textbf{Characterisation} \\
\midrule
$\overline{E} = E \cup E'$
  & closure = the set plus all its limit points \\[2pt]
$\overline{E} = E \cup \partial E$
  & closure = the set plus its boundary \\[2pt]
$E$ closed $\iff$ $E = \overline{E}$
  & closed sets are their own closure \\[2pt]
$\overline{E}$ is the smallest closed set containing $E$
  & closure is the tight closed wrapper around $E$ \\[2pt]
\bottomrule
\end{tabular}
\end{tcolorbox}

\bigskip

% =========================================================
% PART V: IBE DECOMPOSITION — POINT CLASSIFICATION TABLE
% =========================================================

\begin{tcolorbox}[colback=gray!6, colframe=gray!40, arc=2pt,
  left=8pt, right=8pt, top=6pt, bottom=6pt,
  title={\small\textbf{Part V \enspace Point Classification:
         Interior / Boundary / Exterior}},
  fonttitle=\small\bfseries]
\small
Every point $x \in X$ belongs to exactly one region relative to $A \subseteq X$.
\renewcommand{\arraystretch}{1.55}
\begin{tabular}{p{2.2cm} p{4.4cm} p{2.0cm} p{3.6cm}}
\toprule
\textbf{Region} & \textbf{Ball condition} & \textbf{Is $x \in A$?} & \textbf{Is $x \in \overline{A}$?} \\
\midrule

$\operatorname{int}(A)$
  & $\exists r > 0 : B_r(x) \subseteq A$
  & Yes (always)
  & Yes \\[2pt]

$\partial A$
  & $\forall r > 0 : B_r(x)$ meets both $A$ and $X \setminus A$
  & Maybe
  & Yes \\[2pt]

$\operatorname{ext}(A)$
  & $\exists r > 0 : B_r(x) \cap A = \varnothing$
  & No (never)
  & No \\[2pt]

\midrule
\multicolumn{4}{l}{$X = \operatorname{int}(A) \;\dot\cup\; \partial A \;\dot\cup\; \operatorname{ext}(A)$
\quad (disjoint, exhaustive) \quad \hyperref[thm:ibe-decomp]{T\ref*{thm:ibe-decomp}}} \\[2pt]

\midrule
\multicolumn{4}{l}{
  $A$ open $\iff$ $A \cap \partial A = \varnothing$ \qquad
  $A$ closed $\iff$ $\partial A \subseteq A$
} \\[2pt]

\bottomrule
\end{tabular}
\end{tcolorbox}

\bigskip

% =========================================================
% PART VI: TOPOLOGY AXIOMS — OPEN/CLOSED UNDER OPERATIONS
% =========================================================

\begin{tcolorbox}[colback=gray!6, colframe=gray!40, arc=2pt,
  left=8pt, right=8pt, top=6pt, bottom=6pt,
  title={\small\textbf{Part VI \enspace Topology Axioms:
         Operations on Open and Closed Sets}},
  fonttitle=\small\bfseries]
\small
\renewcommand{\arraystretch}{1.55}
\begin{tabular}{p{3.4cm} p{3.6cm} p{4.8cm}}
\toprule
\textbf{Operation} & \textbf{Open sets} & \textbf{Closed sets} \\
\midrule

Arbitrary union
  & $\bigcup_\alpha U_\alpha$ is open
  & (arbitrary intersection is closed) \\[2pt]

Finite intersection
  & $U_1 \cap \cdots \cap U_n$ is open
  & (finite union is closed) \\[2pt]

$\varnothing$ and $X$
  & both open
  & both closed (clopen) \\[2pt]

\midrule
\multicolumn{3}{l}{\textbf{Warning:} infinite intersection of open sets need not be open
  \quad e.g.\ $\bigcap_{n=1}^\infty (-1/n, 1/n) = \{0\}$ in $\mathbb{R}$} \\[2pt]
\multicolumn{3}{l}{\textbf{Warning:} infinite union of closed sets need not be closed
  \quad e.g.\ $\bigcup_{n=1}^\infty [1/n, 1] = (0, 1]$ in $\mathbb{R}$} \\[2pt]
\bottomrule
\end{tabular}

\medskip
\noindent Sources:
\hyperref[prf:union-open-sets]{arbitrary unions of open sets};
\hyperref[prf:intersection-open-sets]{finite intersections of open sets};
\hyperref[prf:intersection-closed-set]{arbitrary intersections of closed sets};
\hyperref[prf:finite-union-closed-sets]{finite unions of closed sets}.
\end{tcolorbox}

\bigskip

% =========================================================
% PART VII: THE FOUR POSSIBILITIES
% =========================================================

\begin{tcolorbox}[colback=gray!6, colframe=gray!40, arc=2pt,
  left=8pt, right=8pt, top=6pt, bottom=6pt,
  title={\small\textbf{Part VII \enspace The Four Possibilities}},
  fonttitle=\small\bfseries]
\small
Open and closed are \emph{independent} properties.
A set in a metric space can be any combination:

\renewcommand{\arraystretch}{1.55}
\begin{tabular}{p{2.4cm} p{2.6cm} p{6.4cm}}
\toprule
\textbf{Open?} & \textbf{Closed?} & \textbf{Example in $\mathbb{R}$} \\
\midrule

Yes & Yes   & $\varnothing$, $\mathbb{R}$ (clopen sets) \\[2pt]
Yes & No    & $(0,1)$, any open interval \\[2pt]
No  & Yes   & $[0,1]$, $\mathbb{Z}$, any closed interval \\[2pt]
No  & No    & $[0,1)$, $\mathbb{Q}$ \\[2pt]

\bottomrule
\end{tabular}

\medskip
\noindent\textbf{Key point:} ``not open'' does not mean closed, and
``not closed'' does not mean open.
The word ``clopen'' names sets that are both.
\end{tcolorbox}
